\section{Discussion}\label{sec:discussion}

\paragraph{Where the error came from.} Before this study we assumed the weak point was sub-pixel accuracy. It is
not. On clean frames every method is within 0.03--0.15\,mm. The baseline's large stress errors come from three
specific faults, each with a specific classical fix: (i) no notion of ``one line'' (fixed by DP tracking), (ii) a
depth stage that fits in image columns and lets one stray point set the floor (fixed by fitting in millimetres and
clustering the floor), and (iii) no ability to say ``retake'' (fixed by the gate).

\paragraph{Classical or learned?} Table~\ref{tab:verdicts} summarises each family. The pattern is consistent:
learned methods looked best on the simulated sets that resemble their training data and lost most when the
pictures changed (another unit, a hand-held photo). The classical methods moved least. Learning earned its place in
one role: the small logistic-regression gates, where the question is fuzzy (``does this look like a usable photo?'')
and a wrong answer costs a retake rather than a wrong measurement.

\begin{table}[H]
\centering\small
\caption{Summary verdict per family.}\label{tab:verdicts}
\begin{tabularx}{\linewidth}{@{}lllX@{}}
\toprule
Family & Simulation & Real photos & Verdict \\
\midrule
DP line tracking & strong & agrees ($\approx$1\,px) & \textbf{Adopt} as the line extractor. \\
Robust depth stage & strong & consistent on the repeat pair & \textbf{Adopt}. Confirm against a gauge. \\
Quality gate & strong & 0.24 false-accept, cross-validated & \textbf{Adopt}; retrain as labelled sessions grow. \\
Calibration (physical model) & strong & not testable yet & \textbf{Adopt} the eight-height recipe and zero check. \\
Laser on/off & best of all & needs hardware & \textbf{Hardware request}: laser on a GPIO. \\
Per-column peaks & accurate, fragile & mixed & Keep as the fast option if the MCU budget is tight. \\
U-Net line & strong & agrees ($\approx$1\,px) & Park until real holder images exist. \\
Candidate classifier & strong & fails & Do not use: learned a simulator shortcut. \\
End-to-end CNN & good in-domain & fails & Do not use as the measurement. \\
\bottomrule
\end{tabularx}
\end{table}

\paragraph{Computing cost.} On one laptop-class core the recommended stack without the frame gate takes about
35--40\,ms per frame, against about 26\,ms for the baseline; the frame gate's features roughly double that and are
the first thing to optimise. Timings on the shared test machine varied by $\pm30\,\%$ between runs, so treat them as
rough. Nothing in stages 1--3 needs more than NumPy-style array operations, so a port to the STM32H7 or to the
cloud back-end is realistic; the U-Net is not realistic on the microcontroller.

\section{Independent audit}\label{sec:audit}

Before writing this report, a separate reviewer that had written none of the code re-ran the key results from a
clean checkout and looked for mistakes. The full audit is in \texttt{research/audit/AUDIT.md}; Table~\ref{tab:audit}
summarises it.

\begin{table}[H]
\centering\small
\caption{Audit summary.}\label{tab:audit}
\begin{tabularx}{\linewidth}{@{}lX@{}}
\toprule
Check & Outcome \\
\midrule
Reproducibility & Nine algorithms re-scored; every one of 89 metrics per algorithm matched the stored value exactly. \\
Data integrity & Test frames regenerate bit-for-bit from their seeds; the six simulated sets share no scenes. \\
Leakage & No learned model was trained or tuned on a test set. One in-sample number: the frame gate on real photos
          (flagged wherever it appears). \\
Scorer & Behaves correctly on synthetic perfect, empty and always-abstain algorithms. Three quirks found (below). \\
\bottomrule
\end{tabularx}
\end{table}

\paragraph{What the audit changed in this report.}
\begin{itemize}
\item \textbf{Means have wide error bars.} A few gross failures dominate every stress MAE: for DP tracking, seven
      grooves carry 78\,\% of the total error. Table~\ref{tab:boot} gives bootstrap intervals. Both improvements over
      the baseline are real; DP tracking and the robust depth stage \emph{cannot be ranked against each other} on
      mean error. We therefore report medians and gross-error rates next to every mean, and we do not rank the
      better methods by small MAE differences.
\item \textbf{Unsafe-PASS rates rest on few tyres.} Only 10, 17 and 14 tyres are illegal in \texttt{ideal},
      \texttt{stress} and \texttt{shifted}; ``0.000'' there means ``0 of 17''. Counts are given in
      Section~\ref{sec:stack}. Verdict accuracy on those sets is also flattering: always answering PASS would score
      0.915 on \texttt{stress}. Only \texttt{legal} (64 illegal tyres of 150) supports a verdict claim.
\item \textbf{A scoring artefact was removed.} The scorer showed the end-to-end network making no false measurements
      on real photos, because it outputs millimetres and the real-photo scorer only counts pixel outputs. Run
      directly, it reports a depth on 43 of 67 unusable photos. Its rows are left out of the real and hand-held
      tables.
\item \textbf{On/off numbers are for a perfectly still device.} In the benchmark the laser-off frame is the same
      scene with no motion. The effect of motion was studied separately (Figure~\ref{fig:onoff}), and there is no
      real-photo evidence for on/off at all.
\item \textbf{The U-Net was trained once}, with one random seed; its numbers have unknown run-to-run spread.
\end{itemize}

\begin{table}[H]
\centering\small
\caption{Bootstrap over frames (1{,}000 resamples) on \texttt{stress}: depth error in mm with 95\,\% intervals.}\label{tab:boot}
\begin{tabular}{@{}lllrr@{}}
\toprule
Algorithm & MAE [95\,\% interval] & Median [95\,\% interval] & Error carried by worst 5\,\% & MAE without them \\
\midrule
Baseline & 1.52 [1.13, 1.93] & 0.54 [0.32, 0.67] & 48\,\% & 0.83 \\
DP line only & 0.54 [0.25, 0.94] & 0.09 [0.07, 0.11] & 78\,\% & 0.12 \\
Robust depth only & 0.46 [0.34, 0.63] & 0.09 [0.07, 0.11] & 53\,\% & 0.23 \\
\midrule
\multicolumn{5}{@{}l}{Paired differences in MAE: baseline $-$ DP $=0.95$ [0.44, 1.45]; baseline $-$ robust depth $=1.03$ [0.68, 1.45];}\\
\multicolumn{5}{@{}l}{DP $-$ robust depth $=0.08$ [$-0.24$, $+0.48$] (no detectable difference).}\\
\bottomrule
\end{tabular}
\end{table}


\section{Threats to validity}\label{sec:validity}

We list these plainly because they decide how far the numbers above can be trusted.

\begin{enumerate}
\item \textbf{Most numbers come from a simulator.} It models geometry, blur, noise, glints and dirt, but its line is
      thinner than the real one, always sits at the same height in the frame in the standard sets, and its tread
      crown is an exact parabola. Two learned methods exploited such regularities. The \texttt{handheld} and
      \texttt{shifted} sets were added to expose this, but they are still simulations.
\item \textbf{The simulator shares assumptions with some methods.} Its camera is exactly the physical calibration
      model, and its crown is exactly what the robust depth stage fits. Gains in those two places are upper bounds.
\item \textbf{The real photos were not taken with the holder.} They are hand-held, the laser was not fixed to the
      camera, and no groove was measured with a gauge. Their depths are meaningful only in pixels, and the six
      reference depths come from the baseline itself. Six frames cannot rank methods that differ by a pixel.
\item \textbf{Small numbers.} 19 usable photos, 12 with an unusable laser. One photo is 5--8 percentage points.
\item \textbf{In-sample gate numbers.} The frame gate was trained on the same 86 photos; we quote cross-validated
      figures wherever we draw a conclusion. All 19 usable photos come from two sessions, so even those are a
      feasibility result rather than a performance estimate.
\item \textbf{Tuning.} All thresholds were tuned on \texttt{train}. Two choices in the combined stack were made
      after test numbers had been seen (the pixel cap, and the 3\,\% risk target); neither changes a conclusion.
\item \textbf{Scorer limits.} There is no ``at least $x$\,mm'' outcome and no guard band at 1.6\,mm, so a 0.02\,mm
      slip across the limit counts the same as a gross unsafe PASS; verdict accuracy is also inflated by the fact
      that most simulated tyres are legal.
\end{enumerate}

\section{Recommendations and next steps}

\paragraph{Software (ready now).}
\begin{enumerate}
\item Replace the Steger line extractor with \algo{dp\_dense} and the depth stage with \algo{rd\_final}; add the
      gates. This is \algo{stack\_dp\_rd\_gate} on branch \texttt{research/best-stack}.
\item Report three outcomes, not two: PASS, FAIL, and \emph{retake}. Add a guard band around 1.6\,mm.
\item Switch calibration to the physical model with eight plate heights from $+4$ to $-18$\,mm, and add a
      one-photo zero check at the start of a shift.
\end{enumerate}

\paragraph{Hardware (one request).} Let the microcontroller switch the laser, and capture an off frame within a
few milliseconds of the on frame. In simulation this is worth more than any other single change.

\paragraph{Data (the real gap).} Everything above should be confirmed on images from the printed holder:
\begin{enumerate}
\item With the holder, photograph 5--10 tyres, three presses per groove, and measure each groove with the manual
      gauge. Fifty gauge readings would turn every ``agreement'' number in this report into an accuracy number.
\item Photograph a machined step block (for example 2, 5 and 10\,mm steps) to check the calibration end to end.
\item Keep labelling ``usable / not usable'' on every session; the gate improves with sessions, not with model size.
\end{enumerate}

\section{Conclusion}

Ten algorithm families were implemented on separate branches and scored on the same simulated and real data.
The most effective changes were not the most sophisticated ones: a dynamic-programming line tracker, a depth stage
that works in millimetres and takes the groove floor from a cluster, and a small gate that can refuse a photo.
In simulation, together they cut the mean depth error under bad conditions from 1.5\,mm to 0.25\,mm, cut gross errors
from 36\,\% to under 3\,\%, passed none of 41 illegal tyres where the baseline passed 9, and stopped the device from reporting
depths on photos with no laser, while still answering about nine presses in ten. Neural networks that replace the
geometry looked good in simulation and failed on real photos and on other units. The next step is not a better
algorithm but ground truth: gauge-measured grooves photographed through the printed holder.

\appendix
\section{Reproducing the results}

\begin{lstlisting}[language=bash]
git clone https://github.com/Advait-B/urex && cd urex
git checkout research/paper                 # everything merged: simulator, benchmark, all algorithms, this paper
pip install numpy scipy scikit-learn scikit-image opencv-python-headless torch
cd research
python -m bench.run2 algos.stack_dp_rd_gate --out results2/stack_dp_rd_gate.json   # score one algorithm
python -m bench.score_all algos.dp_dense algos.rd_final                             # re-score several
python paper/make_tables.py && (cd paper && pdflatex main.tex)                      # rebuild tables and PDF
\end{lstlisting}

The simulated sets are generated from fixed seeds the first time they are needed (a few minutes) and cached.
Each branch has a write-up in \texttt{research/notes/<topic>.md} with more variants, figures and negative results
than fit here.
